\section{Discussion --- draft v1}
\label{sec:discussion}

\subsection{Where the knowledge base stops scaling}

The switch-indexed knowledge base (D4) is honest about its ceiling,
and we hit the first warning shot ourselves: the \texttt{detail}/\texttt{detail2}
parameterization ran out of arity the first time a diagnostic needed
three runtime values, and the factory works around it by re-deriving
one field (\S6.2). At seventeen kinds the workaround is a documented
oddity; at hundreds it would be a code smell; at Clang's or Rust's
scale it is exactly why those projects hold messages in TableGen and
Fluent respectively \cite{rustc-dev-guide,clang-internals}. The pattern's
load-bearing property --- one place where all prose lives, indexed by
root cause --- survives that migration; what changes is the storage
(data files instead of a switch) and the parameterization (named
placeholders instead of positional details). We read our capacity
limit as evidence \emph{for} the industrial designs, reached from the
opposite direction.

\subsection{Trace depth in bigger languages}

Our recorded traces are short because the language is flat: no
recursion in the grammar beyond expressions, no template
instantiation, no macro expansion. In a production front end the
active-stack snapshot could be hundreds of frames deep (consider a
C++ template error's instantiation stack --- which production
compilers already summarize, badly or well, because they face the
same problem). The scope-placement rule (D9) generalizes --- scopes
where diagnostics can fire, plus structural landmarks --- but a
production adoption would need a \emph{budget policy} the rule does not
yet have: depth windows, frame sampling, or collapse-repeated-frames
summarization. We flag this as the first question a scaling attempt
must answer, not as a solved problem.

\subsection{Provenance's price in context}

The 12--17\% always-on cost (\S9.2) looks alarming until placed against
practice: production compilers already carry per-instruction source
locations when building with debug info, and LLVM's remark
infrastructure \cite{llvm-remarks} attaches transformation records opt-in. Our
number is what \emph{always-on, value-carried} provenance costs in a
compiler whose instructions were previously 116 bytes of pure
computation --- a worst-case framing, since any compiler with existing
debug-location machinery has already paid most of the memory-layout
bill. The measured reduction paths are concrete: packing the span
(20 bytes of four ints could be an interned 4-byte index), interning
note strings, or gating note construction the way recording gates
detail construction (D14 --- the same fix, one field over). We did not
apply them because 12--17\% at teaching scale does not obstruct the
thesis, and an optimization pass without a driving requirement is
how measurement artifacts get built.

\subsection{Two observability modes, one architecture}

The implementation carries both a \emph{narrating} channel (the semantic
check log: one line per check performed, always on, educational) and
an \emph{evidentiary} channel (recorded traces: silent until a diagnostic
fires). They answer different questions --- "what is the compiler
doing?" versus "why did THIS happen?" --- and their coexistence is not
redundancy but a division the architecture makes cheap: both are
consumers of the same stage-owned state. The duality suggests the
general shape of compiler observability: continuous narration for
learning contexts, on-demand evidence for diagnosis, one contract
underneath.

\subsection{Retrofit versus designed-in}

GHC's and Rust's structured-diagnostics migrations \cite{ghc-structured-errors,welltyped2021ghc,rust-rfc1644} are the
strongest external evidence this architecture addresses a real need
--- and the strongest evidence about cost asymmetry. Their retrofits
are multi-year, multi-release migrations across entrenched
string-based call sites \cite{ghc-structured-errors,welltyped2021ghc}; our designed-in version cost $\approx$320
lines and 5.6\% of a (small) codebase (\S8.2). The honest generalization is not "designing
in is cheap" --- our codebase was shaped to receive it --- but
narrower: the \emph{marginal} cost of the explanation contract is small
when the failure-handling architecture is decided early, and grows
with every string-typed diagnostic call site added before the
decision.

\subsection{Evidence for the LLM era}

Generated error explanations are being deployed now, with mixed
results \cite{taylor2024dcc,santos2024silver}. Whatever position one takes on wording generation,
an LLM explainer needs \emph{grounding} --- the actual construct, the
actual path, the actual transformation history --- precisely the
evidence conventional pipelines discard and this architecture
preserves as structured values. We speculate, and mark it as
speculation, that recorded traces and provenance form a natural
retrieval substrate for grounded generation; nothing in this paper
tests that.



